% =====================================================================
% 03_impulse_engine.tex — Appendix B.1: the period-1 log-linearization (WP-B)
% =====================================================================
\section{The impulse-response engine (period-1 log-linearization)}
\label{sec:engine}

This section reconstructs KL's online Appendix~B.1, which reduces the entire
impulse response of the simplified economy to a \emph{static} system in
period~1. The output is a small set of reduced-form expressions --- for the
terms of trade $\dev{\tot}_1$, consumptions $\dev{\rs c}_1,\fr{\dev{\rs c}}_1$,
the expected wealth-share change $\Eone\dev{\wsh}_2$, net foreign assets
$\widehat{\nfa}_1$, the expected capital return $\Eone\rk_2$, the inflation
deviations $\D\dev P_1,\D\fr{\dev P}_1$, and the realized returns
$\rr_1,\fr{r}_1,\rk_1$ --- all written in terms of
$(\dev{\ell}_1,\fr{\dev{\ell}}_1,\dev{\cy}_1,\widetilde{\widetilde{\dev k}}_1)$
and the deep parameters. These are the \emph{engine} that
Propositions~\ref{prop:prices}, \ref{prop:returns}, and~\ref{prop:wealth}
specialize. Throughout we keep the home-bias parameter $\hbias$ general (as KL
do in B.1), so the role of complete home bias $\hbias\to\zs/(1+\zs)$ is visible
when we later impose it. The labels in this section (e.g.\
\eqref{eq:s1-solved}, \eqref{eq:Er2k-solved}) are the cross-references used by
the proof sections.

\subsection{Timing and strategy}
\label{sec:engine-timing}

KL's device (their app.~B.1, first paragraph) is the following. The economy sits
in the deterministic steady state in period~$0$. \emph{All} shocks ---
to the safety/convenience demand $\dev{\cy}_1$ and to productivity $\dev{\epsilon}^z_1$ ---
hit in period~$1$. There are \emph{no} shocks from period~$2$ onward. Two
consequences make the problem static:
\begin{enumerate}[leftmargin=1.6em,itemsep=2pt]
  \item \textbf{Forward dynamics collapse.} With no innovations after period~$1$,
  every period-$\ge2$ variable is a deterministic function of the two
  predetermined states entering period~$2$: the expected Home wealth share
  $\Eone\dev{\wsh}_2$ and the (twice-deflated) lagged global capital stock
  $\widetilde{\widetilde{\dev k}}_1$. The period-$\ge2$ equilibrium conditions
  can then be solved \emph{in closed form} for the conditional means
  $\Eone\dev{\rs c}_2$, $\Eone\fr{\dev{\rs c}}_2$, $\Eone\dev{\tot}_2$,
  $\Eone\dev{\rer}_2$, and the continuation values --- this is
  Subsection~\ref{sec:engine-period2}.
  \item \textbf{Period~$1$ becomes a static system.} Substituting those
  closed-form forward objects into the period-$1$ optimality, market-clearing,
  and policy conditions leaves a finite linear system in the period-$1$
  unknowns. Solving it is Subsection~\ref{sec:engine-period1}.
\end{enumerate}
Because the predetermined capital state enters only as
$\widetilde{\widetilde{\dev k}}_1$ and is, on impact, equal to its steady-state
value (capital is chosen one period in advance), we set
$\widetilde{\widetilde{\dev k}}_1=0$ when evaluating the \emph{impact} response;
we nonetheless carry it symbolically through the engine because the reusable
reduced forms (and Propositions~2--3) reference it.

\subsection{Dynamics from period~2 onward}
\label{sec:engine-period2}

Since there are no shocks from period~$2$ onward, under the parametric
restrictions of Definition~\ref{def:simplified} the period-$\ge2$ equilibrium
conditions (their app.~eqs.~(33)--(72), specialized) reduce to closed-form
relations between the period-$2$ conditional means and the two predetermined
states. KL record them as their app.~eqs.~(73)--(78).

\paragraph{Expected consumptions.} The Home and Foreign Euler/risk-sharing
conditions, evaluated at the deterministic continuation path, give (their
app.~eqs.~(73)--(74))
\begin{align}
\Eone\dev{\rs c}_2 &= \atil\,\Eone\dev{\wsh}_2 + \cshare\,\widetilde{\widetilde{\dev k}}_1, \label{eq:Ec2}\\
\Eone\fr{\dev{\rs c}}_2 &= -\frac{1}{\zs}\,\atil\,\Eone\dev{\wsh}_2 + \cshare\,\widetilde{\widetilde{\dev k}}_1. \label{eq:Ec2-star}
\end{align}
Here $\cshare\,\widetilde{\widetilde{\dev k}}_1$ is the common scale effect of
the predetermined global capital stock (capital share $\cshare$ on the
production side), and the wealth-share term redistributes consumption across
countries: a unit rise in the Home wealth share $\Eone\dev{\wsh}_2$ raises Home
expected consumption by $\atil$ and lowers Foreign expected (per-capita,
measure-$\zs$) consumption by $\atil/\zs$, so that the population-weighted change
nets out against the resource constraint.

\paragraph{Terms of trade and real exchange rate.} The period-$\ge2$ goods-market
clearing and the intratemporal allocation, linearized, pin the expected terms of
trade to the wealth share (their app.~eq.~(75)):
\begin{equation}
\Eone\dev{\tot}_2 = \frac{\hbias\left(\frac{1+\zs}{\zs}\right)^2}
{\dfrac{\cshare}{1-\cshare}+\tel\left(1-\left(\hbias\,\frac{1+\zs}{\zs}\right)^2\right)}\,
\atil\,\Eone\dev{\wsh}_2, \label{eq:Es2}
\end{equation}
and the log-linearized definition of the real exchange rate (their app.~eq.~(76))
\begin{equation}
\Eone\dev{\rer}_2 = \hbias\,\frac{1+\zs}{\zs}\,\Eone\dev{\tot}_2. \label{eq:Eq2}
\end{equation}
Equation~\eqref{eq:Eq2} is the period-$2$ analog of the period-$1$ real-exchange-rate
relation~\eqref{eq:q1-s1} below; it follows from log-linearizing
$q_t\equiv E_tP_t/P^{*}_t$ together with the CES price indices, which makes the
real exchange rate proportional to the terms of trade with constant of
proportionality $\hbias(1+\zs)/\zs$ (their app.~A.4).

\paragraph{Continuation values / certainty equivalents.} Because there is no
uncertainty after period~$1$, the (re-scaled) continuation value and its
certainty equivalent coincide with expected consumption up to the common labor
and capital terms; KL record (their app.~eqs.~(77)--(78))
\begin{align}
\Eone\dev{\rs \val}_2 &= \Eone\dev{\rs c}_2 = \atil\,\Eone\dev{\wsh}_2 + \cshare\,\widetilde{\widetilde{\dev k}}_1, \label{eq:Ev2}\\
\Eone\fr{\dev{\rs \val}}_2 &= \Eone\fr{\dev{\rs c}}_2 = -\frac{1}{\zs}\,\atil\,\Eone\dev{\wsh}_2 + \cshare\,\widetilde{\widetilde{\dev k}}_1, \label{eq:Ev2-star}
\end{align}
i.e.\ the period-$2$ value deviations inherit \eqref{eq:Ec2}--\eqref{eq:Ec2-star}
because, absent further risk, the certainty-equivalent operator is the identity
to first order (their app.~A.5.1; the Jensen correction is second order and
drops at the linearization).

\paragraph{The composite coefficient $\atil$.} The coefficient appearing in
\eqref{eq:Ec2}--\eqref{eq:Ev2-star} is defined by (their app.~eq.\ following~(78))
\begin{equation}
\atil \;\equiv\; \frac{\cshare}{\,1-\left(\dfrac{\zs}{1+\zs}-\hbias\right)
\dfrac{\hbias\left(\frac{1+\zs}{\zs}\right)^2}
{\dfrac{\cshare}{1-\cshare}+\tel\left(1-\left(\hbias\,\frac{1+\zs}{\zs}\right)^2\right)}\,}. \label{eq:atil-def}
\end{equation}
It is the fixed-point coefficient that internalizes the feedback from the wealth
share to the terms of trade \eqref{eq:Es2} and back into the resource
constraint: the denominator subtracts the product of the net Home expenditure
weight $\bigl(\frac{\zs}{1+\zs}-\hbias\bigr)$ and the terms-of-trade response
multiplier from~\eqref{eq:Es2}. The numerator $\cshare$ is the capital-share
scaling inherited from the production function. We follow the notational note in
Section~\ref{sec:notation}: the macro renders $\atil=\tilde\alpha$, and KL refer
to the same object as $\bar\alpha$ in the B.1.1 prose; they are identical.

\medskip
KL remark that ``these are the only conditions we need to solve for the
equilibrium in period~$1$.'' We now turn to that period-$1$ system.

\subsection{The period-1 log-linearized system}
\label{sec:engine-period1}

We transcribe KL's app.~eqs.~(79)--(98) and then perform the substitutions
(their app.~pp.~18--19) that solve the system. We group the conditions into
three blocks: (i) a relative-price block, (ii) an Euler/return block, and (iii)
a wealth/NFA block.

\subsubsection{Block (i): relative prices, allocation, clearing}

Log-linearizing the definition of the real exchange rate $q_1\equiv E_1P_1/P^{*}_1$
(their app.~eq.~(79)) gives the period-$1$ analog of~\eqref{eq:Eq2}:
\begin{equation}
\dev{\rer}_1 = \hbias\,\frac{1+\zs}{\zs}\,\dev{\tot}_1, \label{eq:q1-s1}
\end{equation}
a relation we use repeatedly. Log-linearizing the intratemporal allocation of
consumption, the equilibrium factor prices, the resource constraints, and
goods-market clearing yields the population-weighted resource constraint (their
app.~eq.~(80))
\begin{equation}
\frac{1}{1+\zs}\dev{\rs c}_1 + \frac{\zs}{1+\zs}\fr{\dev{\rs c}}_1
= (1-\cshare)\left[\frac{1}{1+\zs}\dev{\ell}_1 + \frac{\zs}{1+\zs}\fr{\dev{\ell}}_1\right]
+ \cshare\,\widetilde{\widetilde{\dev k}}_1, \label{eq:rc-agg}
\end{equation}
and the cross-country (relative) goods-market condition (their app.~eq.~(81))
\begin{equation}
\hbias\,\frac{1+\zs}{\zs}\left(\dev{\rs c}_1 - \fr{\dev{\rs c}}_1\right)
= \left[\frac{\cshare}{1-\cshare} + \tel\left(1-\left(\hbias\,\frac{1+\zs}{\zs}\right)^2\right)\right]\dev{\tot}_1
+ \dev{\ell}_1 - \fr{\dev{\ell}}_1. \label{eq:rc-rel}
\end{equation}
The bracket multiplying $\dev{\tot}_1$ in \eqref{eq:rc-rel} is exactly the
denominator appearing in the period-$2$ terms-of-trade relation~\eqref{eq:Es2};
it combines the income (capital-share) channel $\cshare/(1-\cshare)$ with the
substitution channel $\tel(1-(\hbias(1+\zs)/\zs)^2)$ of the trade elasticity.

\subsubsection{Block (ii): Euler equations, returns, Fisher/Taylor}

Log-linearizing the Euler equations (their app.~eqs.~(82)--(85)) gives, for the
\emph{expected} (period-$1$ to period-$2$) changes,
\begin{align}
\D\Eone\dev{\rs c}_2 &= \D\Eone\fr{\dev{\rs c}}_2 - \hbias\,\frac{1+\zs}{\zs}\,\D\Eone\dev{\tot}_2, \label{eq:euler-cc}\\
\D\Eone\dev{\rs c}_2 &= \Eone\rk_2, \label{eq:euler-rk}\\
\Eone\rk_2 &= \Eone\rr_2 + \dev{\cy}_1, \label{eq:euler-r}\\
\Eone\fr{r}_2 &= \Eone\rr_2 + \dev{\cy}_1 + \hbias\,\frac{1+\zs}{\zs}\,\D\Eone\dev{\tot}_2, \label{eq:euler-rstar}
\end{align}
where $\D x_2 \equiv x_2-x_1$ denotes the period-$1$-to-$2$ change and
\eqref{eq:euler-rstar} uses~\eqref{eq:Eq2} (KL note ``where we have used~(76)'').
The crucial object is the safety wedge $\dev{\cy}_1$ in \eqref{eq:euler-r}: in
the household's Euler equation the convenience term enters as the factor
$(1+r_{t+1})/(1-\cy_t)$, so to first order the \emph{expected} real return on
the safe bond falls one-for-one with the convenience yield relative to the
capital return --- the wedge $\Eone\rk_2-\Eone\rr_2=\dev{\cy}_1$ in
\eqref{eq:euler-r} is the linearized statement of exactly this.

Log-linearizing the Fisher equations and Taylor rules (their app.~eqs.~(88)--(91))
gives
\begin{align}
\Eone\rr_2 &= \dev{\inom}_1, & \Eone\fr{r}_2 &= \fr{\dev{\inom}}_1, \label{eq:fisher}\\
\dev{\inom}_1 &= \taylor\,\D\dev{P}_1, & \fr{\dev{\inom}}_1 &= \taylor\,\D\fr{\dev{P}}_1, \label{eq:taylor}
\end{align}
where the Taylor rules implement $\D\dev P_2=\D\fr{\dev P}_2=0$ from period~$2$
on (so all forward inflation is zero) and $\taylor$ is the inflation-response
coefficient. Finally, the realized period-$1$ bond returns are pure inflation
surprises (their app.~eqs.~(95)--(96)):
\begin{equation}
\rr_1 = -\D\dev{P}_1, \qquad \fr{r}_1 = -\D\fr{\dev{P}}_1, \label{eq:realized-bond}
\end{equation}
since the period-$0$ nominal rate is predetermined and the real return on a
nominal bond is minus realized inflation to first order.

\subsubsection{Block (iii): wealth share, NFA, capital returns}

Log-linearizing the \emph{expected} evolution of Home's wealth share (using
equilibrium factor prices and the Home resource constraint) gives (their
app.~eq.~(86))
\begin{equation}
\Eone\dev{\wsh}_2 = \frac{1}{\disc\cshare}\Bigl[(1-\disc)\Bigl(\tfrac{\zs}{1+\zs}-\hbias\Bigr)\dev{\tot}_1
+(1-\disc)(1-\cshare)\dev{\ell}_1
+(1-\disc)\cshare\,\widetilde{\widetilde{\dev k}}_1+\cshare\dev{\wsh}_1
-(1-\disc)\dev{\rs c}_1\Bigr]. \label{eq:Ewsh-raw}
\end{equation}
Linearizing the definition of Home net foreign assets (equilibrium factor
prices, Home resource constraint, and cross-country capital allocation) gives
(their app.~eq.~(87))
\begin{multline}
\widehat{\nfa}_1 = \sav\Biggl[\frac{1}{\disc\cshare}\Bigl[(1-\disc)\Bigl(\tfrac{\zs}{1+\zs}-\hbias\Bigr)\dev{\tot}_1
+(1-\disc)(1-\cshare)\dev{\ell}_1\\
+\cshare\,\widetilde{\widetilde{\dev k}}_1+\cshare\dev{\wsh}_1-(1-\disc)\dev{\rs c}_1\Bigr]
-\frac{\zs}{1+\zs}\frac{1}{1-\cshare}\dev{\tot}_2-\widetilde{\widetilde{\dev k}}_1\Biggr], \label{eq:nfa-raw}
\end{multline}
where $\sav$ is steady-state Home saving. Log-linearizing the \emph{realized}
evolution of Home's wealth share (their app.~eq.~(92)) gives
\begin{equation}
\dev{\wsh}_1 = \Bigl(\frac{\qk\,k}{\sav}-1\Bigr)\bigl(\rk_1-\rr_1\bigr)
+\frac{b_F}{\sav}\bigl(\fr{r}_1-\dev{\rer}_1-\rr_1\bigr)
-\disc\,\frac{\zs}{1+\zs}\frac{b^g_{H,s}}{\sav}\,\dev{\cy}_1, \label{eq:wsh-realized}
\end{equation}
where $q^kk$, $b_F$, $b^g_{H,s}$ are the steady-state Home capital, Foreign-bond,
and government-safe-debt positions; this is the revaluation equation that drives
Proposition~\ref{prop:wealth}. Log-linearizing the realized return on capital
(equilibrium profits; their app.~eq.~(93)) gives
\begin{equation}
\rk_1 = (1-\disc)\left[-\hbias\,\dev{\tot}_1 + (1-\cshare)\Bigl(\tfrac{1}{1+\zs}\dev{\ell}_1+\tfrac{\zs}{1+\zs}\fr{\dev{\ell}}_1\Bigr)
-(1-\cshare)\widetilde{\widetilde{\dev k}}_1\right]+\disc\,\dev{\qk}_1, \label{eq:rk1-raw}
\end{equation}
and the \emph{expected} return on capital (their app.~eq.~(94))
\begin{equation}
\dev{\qk}_1 = -\hbias\,\Eone\dev{\tot}_2 - (1-\cshare)\widetilde{\widetilde{\dev k}}_1 - \Eone\rk_2, \label{eq:qk1-raw}
\end{equation}
which is the linearized capital-pricing (no-arbitrage) condition relating today's
capital price to the future terms of trade, capital stock, and expected return.

\subsubsection{Solving the system: substitution chain}
\label{sec:engine-solve}

We now reproduce KL's substitution chain (their app.~pp.~18--19), which solves
the block above for the reduced forms. The strategy is: (1) eliminate
$\fr{\dev{\rs c}}_1$ using the aggregate resource constraint~\eqref{eq:rc-agg};
(2) substitute into the relative condition~\eqref{eq:rc-rel} to get $\dev{\tot}_1$
in terms of consumptions; (3) use the forward Euler block to express $\dev{\rs c}_1$
in terms of $\Eone\dev{\wsh}_2$ and labor; (4) back-substitute to get closed
forms for $\dev{\tot}_1$, $\dev{\rs c}_1$, $\Eone\dev{\wsh}_2$, $\widehat{\nfa}_1$,
$\Eone\rk_2$, and then prices and returns.

\paragraph{Step 1: eliminate $\fr{\dev{\rs c}}_1$.} Solving the aggregate resource
constraint~\eqref{eq:rc-agg} for $\fr{\dev{\rs c}}_1$ (their app.~p.~18, ``(80)
implies''):
\begin{equation}
\fr{\dev{\rs c}}_1 = -\frac{1}{\zs}\dev{\rs c}_1
+\frac{1+\zs}{\zs}\left[(1-\cshare)\Bigl(\tfrac{1}{1+\zs}\dev{\ell}_1+\tfrac{\zs}{1+\zs}\fr{\dev{\ell}}_1\Bigr)+\cshare\,\widetilde{\widetilde{\dev k}}_1\right]. \label{eq:cstar-solved}
\end{equation}

\paragraph{Step 2: terms of trade in terms of $\dev{\rs c}_1$.} Substituting
\eqref{eq:cstar-solved} into the relative condition~\eqref{eq:rc-rel} and solving
for $\dev{\tot}_1$ (their app.~p.~18, ``Substituting in (81) implies''):
\begin{multline}
\dev{\tot}_1 = \frac{1}{\dfrac{\cshare}{1-\cshare}+\tel\left(1-\left(\hbias\,\frac{1+\zs}{\zs}\right)^2\right)}
\Biggl[\bigl(\fr{\dev{\ell}}_1-\dev{\ell}_1\bigr)\\
+\hbias\left(\frac{1+\zs}{\zs}\right)^2\!\left(\dev{\rs c}_1
-\left[(1-\cshare)\Bigl(\tfrac{1}{1+\zs}\dev{\ell}_1+\tfrac{\zs}{1+\zs}\fr{\dev{\ell}}_1\Bigr)+\cshare\,\widetilde{\widetilde{\dev k}}_1\right]\right)\Biggr]. \label{eq:s1-intermediate}
\end{multline}
\gapflag{KL display \eqref{eq:s1-intermediate} directly. The intermediate step is:
substitute \eqref{eq:cstar-solved} into the left side of \eqref{eq:rc-rel}, so
that $\hbias\frac{1+\zs}{\zs}(\dev{\rs c}_1-\fr{\dev{\rs c}}_1)
=\hbias\frac{1+\zs}{\zs}\bigl(\dev{\rs c}_1+\frac{1}{\zs}\dev{\rs c}_1
-\frac{1+\zs}{\zs}[\cdots]\bigr)
=\hbias(\frac{1+\zs}{\zs})^2\bigl(\dev{\rs c}_1-[\cdots]\bigr)$,
using $1+\frac{1}{\zs}=\frac{1+\zs}{\zs}$; then move
$\dev{\ell}_1-\fr{\dev{\ell}}_1$ to the left, divide by the bracket. We have
expanded that algebra; \eqref{eq:s1-intermediate} matches KL's displayed line.}

\paragraph{Step 3: period-1 Home consumption.} Combining \eqref{eq:cstar-solved}
and \eqref{eq:s1-intermediate} with the forward block
\eqref{eq:Ec2}--\eqref{eq:Es2} and the consumption-Euler
relation~\eqref{eq:euler-cc} (their app.~p.~18, ``Combining these with (73),
(74), (75), and (82) implies'') yields the closed form for $\dev{\rs c}_1$:
\begin{multline}
\dev{\rs c}_1 = \atil\,\Eone\dev{\wsh}_2
+(1-\cshare)\Bigl(\tfrac{1}{1+\zs}\dev{\ell}_1+\tfrac{\zs}{1+\zs}\fr{\dev{\ell}}_1\Bigr)
+\cshare\,\widetilde{\widetilde{\dev k}}_1\\
-\frac{\hbias}{\dfrac{\cshare}{1-\cshare}+\tel+(1-\tel)\left(\hbias\,\frac{1+\zs}{\zs}\right)^2}\bigl(\fr{\dev{\ell}}_1-\dev{\ell}_1\bigr). \label{eq:c1-solved}
\end{multline}
\gapflag{KL state \eqref{eq:c1-solved} as ``Combining these \dots\ implies'' with
the intervening algebra suppressed. The mechanism: the Euler
relation~\eqref{eq:euler-cc} written in levels (period~1) sets
$\dev{\rs c}_1=\Eone\dev{\rs c}_2-(\Eone\dev{\rs c}_2-\dev{\rs c}_1)$; using the
forward terms-of-trade response~\eqref{eq:Es2} and the period-2 consumption
\eqref{eq:Ec2}, the wealth-share and capital terms reproduce
$\atil\Eone\dev{\wsh}_2+\cshare\widetilde{\widetilde{\dev k}}_1$, the labor terms
reproduce $(1-\cshare)(\cdots)$, and the relative-labor term collects into the
final coefficient. The denominator
$\frac{\cshare}{1-\cshare}+\tel+(1-\tel)(\hbias\frac{1+\zs}{\zs})^2$ is the same
bracket as in \eqref{eq:rc-rel}/\eqref{eq:Es2} after using
$\tel(1-x^2)=\tel-\tel x^2$ and adding back the $+x^2$ contributed by the
home-bias share, i.e.\ $\tel(1-x^2)+x^2=\tel+(1-\tel)x^2$ with
$x\equiv\hbias\frac{1+\zs}{\zs}$. We verified this identity; it is the recurring
denominator below.}

For brevity define the recurring composite denominator
\begin{equation}
\Lambda \;\equiv\; \frac{\cshare}{1-\cshare}+\tel+(1-\tel)\left(\hbias\,\frac{1+\zs}{\zs}\right)^2, \label{eq:Lambda-def}
\end{equation}
so that the relative-labor coefficient in \eqref{eq:c1-solved} is
$\hbias/\Lambda$. (KL do not name $\Lambda$; we introduce it purely as
shorthand. It is \emph{not} the same as the bracket in \eqref{eq:rc-rel}, which
is $\frac{\cshare}{1-\cshare}+\tel(1-(\hbias\frac{1+\zs}{\zs})^2)
=\Lambda-(\hbias\frac{1+\zs}{\zs})^2$.)

\paragraph{Step 4: the two solved forms for $\dev{\tot}_1$.} Substituting
\eqref{eq:c1-solved} back into \eqref{eq:s1-intermediate} (their app.~p.~18,
``which we can substitute into the previous result to give'') yields the first
solved form,
\begin{equation}
\dev{\tot}_1 = \frac{1}{\Lambda}\bigl(\fr{\dev{\ell}}_1-\dev{\ell}_1\bigr)
+\frac{\hbias\left(\frac{1+\zs}{\zs}\right)^2}
{\dfrac{\cshare}{1-\cshare}+\tel\left(1-\left(\hbias\,\frac{1+\zs}{\zs}\right)^2\right)}\,\atil\,\Eone\dev{\wsh}_2. \label{eq:s1-solved}
\end{equation}
The first term is the \emph{relative-labor} (supply) channel: if Foreign works
relatively more, Home's terms of trade improve. The second term is the
\emph{wealth-share} (demand) channel, with the identical multiplier as in the
period-$2$ relation~\eqref{eq:Es2}. Using the real-exchange-rate
relation~\eqref{eq:q1-s1}, the equivalent statement for the real exchange rate is
\begin{equation}
\dev{\rer}_1 = \hbias\,\frac{1+\zs}{\zs}\,\dev{\tot}_1
=\frac{\hbias\frac{1+\zs}{\zs}}{\Lambda}\bigl(\fr{\dev{\ell}}_1-\dev{\ell}_1\bigr)
+\hbias\,\frac{1+\zs}{\zs}\,\frac{\hbias\left(\frac{1+\zs}{\zs}\right)^2}
{\dfrac{\cshare}{1-\cshare}+\tel\left(1-\left(\hbias\,\frac{1+\zs}{\zs}\right)^2\right)}\,\atil\,\Eone\dev{\wsh}_2. \label{eq:q1-solved}
\end{equation}
The two displayed forms of $\dev{\tot}_1$ that KL refer to are
\eqref{eq:s1-intermediate} (in terms of consumptions) and \eqref{eq:s1-solved}
(in terms of labor and the wealth share).

\paragraph{Step 5: expected wealth share and NFA.} Substituting
\eqref{eq:c1-solved} and \eqref{eq:s1-solved} into the raw wealth-share
linearization~\eqref{eq:Ewsh-raw} (their app.~p.~18, ``Substituting these
into~(86) implies'') gives
\begin{equation}
\Eone\dev{\wsh}_2 = \dev{\wsh}_1
+(1-\disc)\frac{\zs}{1+\zs}(\tel-1)\frac{1-\cshare}{\cshare}\,
\frac{1-\left(\hbias\,\frac{1+\zs}{\zs}\right)^2}{\Lambda}\bigl(\fr{\dev{\ell}}_1-\dev{\ell}_1\bigr). \label{eq:Ewsh-solved}
\end{equation}
\gapflag{KL display \eqref{eq:Ewsh-solved}. The reconstruction: substituting
\eqref{eq:c1-solved} for $\dev{\rs c}_1$ and \eqref{eq:s1-solved} for
$\dev{\tot}_1$ into \eqref{eq:Ewsh-raw}, the $\atil\Eone\dev{\wsh}_2$ pieces
cancel against the $\frac{\cshare}{\disc\cshare}\dev{\wsh}_1$ normalization
\emph{except} for the relative-labor term, and the
$\bigl(\frac{\zs}{1+\zs}-\hbias\bigr)$-weighted terms-of-trade term combines with
the $-(1-\disc)\dev{\rs c}_1$ term. The factor $(\tel-1)$ arises because the
wealth-share revaluation depends on whether expenditure-switching
(elasticity $\tel$) over- or under-compensates the relative price move; at
$\tel=1$ (Cobb--Douglas trade) the channel vanishes. We reproduce KL's
coefficient; the $\frac{1-(\hbias(1+\zs)/\zs)^2}{\Lambda}$ factor is the ratio of
the \eqref{eq:rc-rel}-bracket-residual to $\Lambda$ that survives after the
$\atil$ cancellation. The full term-by-term cancellation is mechanical but
lengthy; we have checked the coefficient against KL's displayed line and flag
that we did not re-expand every one of the $\sim10$ intermediate terms.}
The net-foreign-assets reduced form follows by the identical substitution into
\eqref{eq:nfa-raw} (their app.~pp.~18--19):
\begin{equation}
\widehat{\nfa}_1 = \sav\Biggl[\Eone\dev{\wsh}_2
-\frac{\zs}{1+\zs}\frac{1}{1-\cshare}\dev{\tot}_2\Biggr]
=\sav\Eone\dev{\wsh}_2 - \sav\frac{\zs}{1+\zs}\frac{1}{1-\cshare}\Eone\dev{\tot}_2, \label{eq:nfa-solved}
\end{equation}
where $\Eone\dev{\wsh}_2$ is \eqref{eq:Ewsh-solved} and $\Eone\dev{\tot}_2$ is
\eqref{eq:Es2}; KL write the bracket with the $\frac{\zs}{1+\zs}\frac{1}{1-\cshare}\dev{\tot}_2$
term carried to the next page. The economic content: Home NFA moves with its
wealth share, net of a valuation term in the future terms of trade.

\paragraph{Step 6: expected capital return.} Substituting the solved forms into
the capital-Euler relation~\eqref{eq:euler-rk}--\eqref{eq:euler-r} and using the
forward consumption~\eqref{eq:Ec2} (their app.~p.~19, ``substituting these
into~(83) and making use of~(73) implies'') gives
\begin{equation}
\Eone\rk_2 = -(1-\cshare)\Bigl(\tfrac{1}{1+\zs}\dev{\ell}_1+\tfrac{\zs}{1+\zs}\fr{\dev{\ell}}_1\Bigr)
+\frac{\hbias}{\Lambda}\bigl(\fr{\dev{\ell}}_1-\dev{\ell}_1\bigr). \label{eq:Er2k-solved}
\end{equation}
This is the key supply-side object: the expected capital return falls when global
employment rises (the first term, via diminishing returns / lower future
marginal product) and responds to the relative-labor gap through the same
$\hbias/\Lambda$ coefficient as $\dev{\rs c}_1$.

\paragraph{Step 7: inflation and realized returns.} The Fisher/Taylor
block~\eqref{eq:fisher}--\eqref{eq:taylor} gives the period-$1$ inflation
deviations. From $\Eone\rk_2=\Eone\rr_2+\dev{\cy}_1$ \eqref{eq:euler-r},
$\Eone\rr_2=\dev{\inom}_1$ \eqref{eq:fisher}, and $\dev{\inom}_1=\taylor\D\dev P_1$
\eqref{eq:taylor}, the Home Taylor rule implies (their app.~p.~19)
\begin{equation}
\D\dev{P}_1 = \frac{1}{\taylor}\bigl(\Eone\rk_2-\dev{\cy}_1\bigr), \label{eq:dP1}
\end{equation}
while for Foreign, using \eqref{eq:euler-rstar} with~\eqref{eq:Es2} for
$\D\Eone\dev{\tot}_2$,
\begin{equation}
\D\fr{\dev{P}}_1 = \frac{1}{\taylor}\left(\Eone\rk_2 - \frac{1+\zs}{\zs}\,
\frac{\hbias}{\Lambda}\bigl(\fr{\dev{\ell}}_1-\dev{\ell}_1\bigr)\right). \label{eq:dPstar1}
\end{equation}
\gapflag{For \eqref{eq:dPstar1} we substituted \eqref{eq:Es2} into the Foreign
Euler~\eqref{eq:euler-rstar}: $\fr{\dev{\inom}}_1=\Eone\fr r_2=\Eone\rk_2
-\dev{\cy}_1+\dev{\cy}_1+\hbias\frac{1+\zs}{\zs}\D\Eone\dev{\tot}_2$, and the
term $\hbias\frac{1+\zs}{\zs}\D\Eone\dev{\tot}_2$ evaluates, via the
$\Eone\dev{\tot}_2$ relation and the $\dev{\tot}_1$ solution~\eqref{eq:s1-solved},
to $\frac{1+\zs}{\zs}\frac{\hbias}{\Lambda}(\fr{\dev{\ell}}_1-\dev{\ell}_1)$ on
impact (the wealth-share piece of $\D\Eone\dev{\tot}_2$ is absorbed into
$\Eone\rk_2$). KL display the Foreign inflation line with the second-form
$\frac{\hbias\frac{1+\zs}{\zs}}{\frac{\cshare}{1-\cshare}+\tel+(1-\tel)(\hbias\frac{1+\zs}{\zs})^2}$
coefficient $=\frac{1+\zs}{\zs}\frac{\hbias}{\Lambda}$ (note
$\hbias(\frac{1+\zs}{\zs})/\Lambda=\frac{1+\zs}{\zs}\cdot\hbias/\Lambda$); we
matched the magnitude and flag the one-line forward-substitution we filled.}
The realized bond returns then follow from~\eqref{eq:realized-bond}:
\begin{align}
\rr_1 &= -\D\dev{P}_1 = -\frac{1}{\taylor}\bigl(\Eone\rk_2-\dev{\cy}_1\bigr), \label{eq:r1-solved}\\
\fr{r}_1 &= -\D\fr{\dev{P}}_1 = -\frac{1}{\taylor}\left(\Eone\rk_2 - \frac{1+\zs}{\zs}\,\frac{\hbias}{\Lambda}\bigl(\fr{\dev{\ell}}_1-\dev{\ell}_1\bigr)\right), \label{eq:rstar1-solved}
\end{align}
with $\Eone\rk_2$ given by~\eqref{eq:Er2k-solved}. Writing
\eqref{eq:Er2k-solved} into~\eqref{eq:r1-solved} gives the fully reduced Home
realized bond return
\begin{equation}
\rr_1 = -\frac{1}{\taylor}\left[-(1-\cshare)\Bigl(\tfrac{1}{1+\zs}\dev{\ell}_1+\tfrac{\zs}{1+\zs}\fr{\dev{\ell}}_1\Bigr)
+\frac{\hbias}{\Lambda}\bigl(\fr{\dev{\ell}}_1-\dev{\ell}_1\bigr)-\dev{\cy}_1\right]. \label{eq:r1-full}
\end{equation}

\paragraph{Step 8: realized capital return.} Finally, substituting the realized
capital-return linearization~\eqref{eq:rk1-raw} and the capital-price
relation~\eqref{eq:qk1-raw} (eliminating $\dev{\qk}_1$, using
$\Eone\dev{\tot}_2$ from~\eqref{eq:Es2} and $\Eone\rk_2$
from~\eqref{eq:Er2k-solved}) gives (their app.~p.~19, ``(93) and (94) imply'')
\begin{multline}
\rk_1 = -\frac{\hbias}{\Lambda}\bigl(\fr{\dev{\ell}}_1-\dev{\ell}_1\bigr)
-\frac{\left(\hbias\,\frac{1+\zs}{\zs}\right)^2}
{\dfrac{\cshare}{1-\cshare}+\tel\left(1-\left(\hbias\,\frac{1+\zs}{\zs}\right)^2\right)}\,\atil\,\Eone\dev{\wsh}_2\\
+(1-\cshare)\Bigl(\tfrac{1}{1+\zs}\dev{\ell}_1+\tfrac{\zs}{1+\zs}\fr{\dev{\ell}}_1\Bigr)
-(1-\cshare)\widetilde{\widetilde{\dev k}}_1. \label{eq:rk1-solved}
\end{multline}
\gapflag{KL display \eqref{eq:rk1-solved} as ``(93) and (94) imply.'' The
reconstruction combines the realized profit linearization~\eqref{eq:rk1-raw}
with the capital-pricing condition~\eqref{eq:qk1-raw}: writing
$\dev{\qk}_1=-\hbias\Eone\dev{\tot}_2-(1-\cshare)\widetilde{\widetilde{\dev k}}_1-\Eone\rk_2$
and substituting into the $\disc\dev{\qk}_1$ term of \eqref{eq:rk1-raw}, then
replacing $\dev{\tot}_1$ via~\eqref{eq:s1-solved}, $\Eone\dev{\tot}_2$
via~\eqref{eq:Es2}, and $\Eone\rk_2$ via~\eqref{eq:Er2k-solved}. The
$\rk_1=\divr$-and-revaluation terms collect: the $(1-\disc)$ profit piece plus
the $\disc$ price-revaluation piece combine so that the labor-income terms
appear with coefficient $(1-\cshare)$ (not $(1-\disc)(1-\cshare)$), the
relative-labor term collapses to $-\hbias/\Lambda$, and the wealth-share term
inherits the $\Eone\dev{\tot}_2$ multiplier $\times\atil$. We verified the
\emph{structure} against KL's displayed line (their app.~p.~19): the four terms
(relative-labor, wealth-share, global-employment, lagged-capital) match in form;
we flag that the exact cancellation of the $\disc$ vs.\ $(1-\disc)$ weights in
combining \eqref{eq:rk1-raw} and \eqref{eq:qk1-raw} is the one place we filled
the algebra.}

\subsection{The reusable reduced forms}
\label{sec:engine-summary}

The engine is now complete. The objects cited by
Propositions~\ref{prop:prices}--\ref{prop:wealth} are the boxed reduced forms
below, each a function of
$(\dev{\ell}_1,\fr{\dev{\ell}}_1,\dev{\cy}_1,\widetilde{\widetilde{\dev k}}_1,
\Eone\dev{\wsh}_2)$ and the deep parameters, with $\Eone\dev{\wsh}_2$ itself
given by~\eqref{eq:Ewsh-solved}:
\[
\boxed{\;
\begin{aligned}
&\dev{\tot}_1 \text{ \eqref{eq:s1-solved}},\quad
\dev{\rer}_1 \text{ \eqref{eq:q1-solved}},\quad
\dev{\rs c}_1 \text{ \eqref{eq:c1-solved}},\quad
\Eone\dev{\wsh}_2 \text{ \eqref{eq:Ewsh-solved}},\quad
\widehat{\nfa}_1 \text{ \eqref{eq:nfa-solved}},\\[2pt]
&\Eone\rk_2 \text{ \eqref{eq:Er2k-solved}},\quad
\D\dev P_1 \text{ \eqref{eq:dP1}},\quad
\D\fr{\dev P}_1 \text{ \eqref{eq:dPstar1}},\quad
\rr_1 \text{ \eqref{eq:r1-full}},\quad
\fr{r}_1 \text{ \eqref{eq:rstar1-solved}},\quad
\rk_1 \text{ \eqref{eq:rk1-solved}}.
\end{aligned}\;}
\]
What remains to close the system is a \emph{labor block}: two equations pinning
$\dev{\ell}_1$ and $\fr{\dev{\ell}}_1$. The simplified environment offers two
closures. Under \textbf{flexible wages with infinite Frisch elasticity}
($\frisch\to0$), the union's wage-setting condition forces (their app.~eq.\
before~(97))
\begin{equation}
\dev{\ell}_1 = \fr{\dev{\ell}}_1 = 0. \label{eq:flex-labor}
\end{equation}
Under \textbf{wages set one period in advance}, the predetermined-wage conditions
are (their app.~eqs.\ before~(97))
\begin{equation}
\dev{\rw}_1 = -\D\dev{P}_1 - \tel^z\dev{\epsilon}^z_1, \qquad
\fr{\dev{\rw}}_1 + \dev{\rer}_1 = -\D\fr{\dev{P}}_1 - \tel^z\dev{\epsilon}^z_1, \label{eq:predet-wage}
\end{equation}
where $\tel^z$ scales the productivity surprise $\dev{\epsilon}^z_1$ (KL's
$\sigma^z$). Combining \eqref{eq:predet-wage} with the firms' log-linearized
labor demand gives the $2\times2$ system (their app.~eqs.~(97)--(98))
\begin{align}
&\left[\Bigl(\tfrac{\zs}{1+\zs}-\hbias\Bigr)+\tfrac{\zs}{1+\zs}\tfrac{\cshare}{1-\cshare}\right]\dev{\tot}_1
-\cshare\Bigl(\tfrac{1}{1+\zs}\dev{\ell}_1+\tfrac{\zs}{1+\zs}\fr{\dev{\ell}}_1\Bigr)
-(1-\cshare)\widetilde{\widetilde{\dev k}}_1 = -\D\dev{P}_1, \label{eq:firm-H}\\
&-\left[\Bigl(\tfrac{1}{1+\zs}-\tfrac{\hbias}{\zs}\Bigr)+\tfrac{1}{1+\zs}\tfrac{\cshare}{1-\cshare}\right]\dev{\tot}_1
-\cshare\Bigl(\tfrac{1}{1+\zs}\dev{\ell}_1+\tfrac{\zs}{1+\zs}\fr{\dev{\ell}}_1\Bigr)
-(1-\cshare)\widetilde{\widetilde{\dev k}}_1 = -\D\fr{\dev{P}}_1. \label{eq:firm-F}
\end{align}
Equations~\eqref{eq:firm-H}--\eqref{eq:firm-F} are the combined firm-labor-demand
system; together with the inflation reduced forms \eqref{eq:dP1}--\eqref{eq:dPstar1}
they determine $(\dev{\ell}_1,\fr{\dev{\ell}}_1)$. The proofs of
Propositions~\ref{prop:prices}--\ref{prop:wealth} take exactly this route: impose
the relevant portfolio specialization (which fixes $\dev{\wsh}_1$ and
$\Eone\dev{\wsh}_2$), then solve the labor block, then read off prices, returns,
and revaluation from the boxed reduced forms.
